\documentclass[11pt]{article}
\usepackage[a4paper,margin=18mm]{geometry}
\usepackage[no-math]{fontspec}
\setmainfont{Latin Modern Roman}
\usepackage{amsmath,amssymb,amsthm,xfrac,xspace,hyperref,graphicx,comment}
\excludecomment{DefTralics}

\providecommand{\bibliofont}{\small}
\newcommand{\ie}{i.e.,\xspace}
\newcommand{\eg}{e.g.,\xspace}
\newcommand{\etc}{etc.\xspace}
\newcommand{\cf}{cf.\xspace}
\newcommand{\resp}{resp.\xspace}
\newcommand{\smaller}{\small}
\newcommand{\Subsection}{\subsection}
\newcommand{\Subsubsection}{\subsubsection}
\newcommand{\skpt}{\smallskip}
\emergencystretch=3em
\numberwithin{equation}{section}\input{author-preamble.tex}
\begin{document}
\begin{center}{\Large Stable item 2049}\end{center}
\paragraph{Source and attribution.}
Alexis Drouot, \emph{Ubiquity of conical points in topological insulators}, Journal de l’École polytechnique — Mathématiques 8 (2021), publisher version, DOI 10.5802/jep.152. \href{https://jep.centre-mersenne.org/articles/10.5802/jep.152/}{Publisher article}. Authors' text: \href{https://creativecommons.org/licenses/by/4.0/}{CC BY 4.0}.
\paragraph{Editorial problem boundary.}
Prove only Theorem 3 below for a smooth compact boundaryless three-manifold and \(N\ge2\) Hermitian matrix families in the printed smooth Fréchet topology. The set of families having only conical degeneracies is open and dense. Both openness and the entire finite-chart global density construction are retained. Physical current asymptotics, Chern-number formulae and differential-operator potential claims are surrounding context only. Original introductory figure artefacts are preserved faithfully rather than interpreted as mathematical corrections.
\paragraph{Difficulty (editorial): H2.}
Sard's theorem for two-by-two matrices does not directly give a global open dense statement for arbitrary Hermitian families on a three-manifold. The author proves stability through the discriminant and local spectral reductions, then recursively perturbs the two-dimensional spectral blocks over a finite chart cover while preserving the already treated regions. That global compatibility construction is the nonroutine obligation behind the stated density result.
\paragraph{Editorial transcription note.}
Original author language, mathematics and ordering are preserved. Publisher front matter is replaced by this independent article wrapper. Bibliography is recovered from matching cached publisher HTML, including its TeX formulas. No new proof or mathematical repair is supplied. Surrounding original results are retained for conservative context and are not extra selected corpus claims.
\clearpage
\input{author-body.tex}
\clearpage
\section*{Editable original Figure 5 text}
\noindent Editorial source note: the following is an exact visual transcription of the authors' prose and mathematical notation embedded in the original vector Figure 5 on physical source page 17. The unchanged original graphic is also retained. The complete corresponding primary proof remains the exact native author text above; no assistant argument or mathematical repair is added.
\input{author-figure5-text.tex}
\begin{thebibliography}{99}
\bibitem{ADH20} H. Ammari, B. Davies \& E. O. Hiltunen - “Robust edge modes in dislocated systems of subwavelength resonators” , 2020
\bibitem{ADH19} H. Ammari, B. Davies, E. O. Hiltunen \& S. Yu - “Topologically protected edge modes in one-dimensional chains of subwavelength resonators” , 2019
\bibitem{AFL18} H. Ammari, B. Fitzpatrick, E. O. Hiltunen, H. Lee \& S. Yu - “Honeycomb-lattice Minnaert bubbles” , 2018
\bibitem{A95} V. I. Arnold - “Remarks on eigenvalues and eigenvectors of Hermitian matrices, Berry phase, adiabatic connections and quantum Hall effect” , Selecta Math. (N.S.) 1 (1995) no. 1, p. 1-19
\bibitem{AS78} J. E. Avron \& B. Simon - “Analytic properties of band functions” , Ann. Physics 110 (1978) no. 1, p. 85-101
\bibitem{ASV13} J. C. Avila, H. Schulz-Baldes \& C. Villegas-Blas - “Topological invariants of edge states for periodic two-dimensional models” , Math. Phys. Anal. Geom. 16 (2013) no. 2, p. 137-170
\bibitem{B19b} G. Bal - “Continuous bulk and interface description of topological insulators” , J. Math. Phys. 60 (2019) no. 8, p. 081506, 20
\bibitem{B19a} G. Bal - “Topological protection of perturbed edge states” , Commun. Math. Sci. 17 (2019) no. 1, p. 193-225
\bibitem{B19c} G. Bal - “Topological invariants for interface modes” , 2019
\bibitem{B84} M. V. Berry - “Quantal phase factors accompanying adiabatic changes” , Proc. Roy. Soc. London Ser. A 392 (1984) no. 1802, p. 45-57
\bibitem{BKR17} C. Bourne, J. Kellendonk \& A. Rennie - “The $K$-theoretic bulk-edge correspondence for topological insulators” , Ann. Inst. H. Poincaré Phys. Théor. 18 (2017) no. 5, p. 1833-1866
\bibitem{C91} Y. Colin de Verdière - “Sur les singularités de van Hove génériques” , in Analyse globale et physique mathématique (Lyon, 1989) , Mém. Soc. Math. France (N.S.) , vol. 46 , Société Mathématique de France, Paris, 1991, p. 99-110
\bibitem{DE99} L. Dieci \& T. Eirola - “On smooth decompositions of matrices” , SIAM J. Matrix Anal. Appl. 20 (1999) no. 3, p. 800-819
\bibitem{D11} M. Domokos - “Discriminant of symmetric matrices as a sum of squares and the orthogonal group” , Comm. Pure Appl. Math. 64 (2011) no. 4, p. 443-465
\bibitem{DP12} L. Dieci \& A. Pugliese - “Hermitian matrices depending on three parameters: coalescing eigenvalues” , Linear Algebra Appl. 436 (2012) no. 11, p. 4120-4142
\bibitem{D19b} A. Drouot - “The bulk-edge correspondence for continuous honeycomb lattices” , Comm. Partial Differential Equations 44 (2019) no. 12, p. 1406-1430
\bibitem{D19a} A. Drouot - “Characterization of edge states in perturbed honeycomb structures” , Pure Appl. Anal. 1 (2019) no. 3, p. 385-445
\bibitem{D19c} A. Drouot - “Microlocal analysis of the bulk-edge correspondence” (2019), arXiv:1909.10474
\bibitem{DW19} A. Drouot \& M. I. Weinstein - “Edge states and the valley Hall effect” , Adv. Math. 368 (2020), p. 107142, 51
\bibitem{EGS05} A. Elgart, G. M. Graf \& J. H. Schenker - “Equality of the bulk and edge Hall conductances in a mobility gap” , Comm. Math. Phys. 259 (2005) no. 1, p. 185-221
\bibitem{FC13} M. Fruchart \& D. Carpentier - “An introduction to topological insulators” , Comptes Rendus Physique 14 (2013) no. 9, p. 779-815
\bibitem{F04} C. Fermanian Kammerer - “Semiclassical analysis of generic codimension 3 crossings” , Internat. Math. Res. Notices (2004) no. 45, p. 2391-2435
\bibitem{FG03} C. Fermanian Kammerer \& P. Gérard - “A Landau-Zener formula for non-degenerated involutive codimension 3 crossings” , Ann. Inst. H. Poincaré Phys. Théor. 4 (2003) no. 3, p. 513-552
\bibitem{FLW16} C. L. Fefferman, J. P. Lee-Thorp \& M. I. Weinstein - “Edge states in honeycomb structures” , Ann. PDE 2 (2016) no. 2, article ID 12, 80 pages
\bibitem{FLW18} C. L. Fefferman, J. P. Lee-Thorp \& M. I. Weinstein - “Honeycomb Schrödinger operators in the strong binding regime” , Comm. Pure Appl. Math. 71 (2018) no. 6, p. 1178-1270
\bibitem{FT16} S. Freund \& S. Teufel - “Peierls substitution for magnetic Bloch bands” , Ann. PDE 9 (2016) no. 4, p. 773-811
\bibitem{FW12} C. L. Fefferman \& M. I. Weinstein - “Honeycomb lattice potentials and Dirac points” , J. Amer. Math. Soc. 25 (2012) no. 4, p. 1169-1220
\bibitem{GP74} V. Guillemin \& A. Pollack - Differential topology , Prentice-Hall, Inc., Englewood Cliffs, N.J., 1974
\bibitem{GP13} G. M. Graf \& M. Porta - “Bulk-edge correspondence for two-dimensional topological insulators” , Comm. Math. Phys. 324 (2013) no. 3, p. 851-895
\bibitem{H93} Y. Hatsugai - “Chern number and edge states in the integer quantum Hall effect” , Phys. Rev. Lett. 71 (1993) no. 22, p. 3697-3700
\bibitem{I92} N. V. Ilyushechkin - “The discriminant of the characteristic polynomial of a normal matrix” , Mat. Zametki 51 (1992) no. 3, p. 16-23, 143
\bibitem{KP07} P. Kuchment \& O. Post - “On the spectra of carbon nano-structures” , Comm. Math. Phys. 275 (2007) no. 3, p. 805-826
\bibitem{KRS02} J. Kellendonk, T. Richter \& H. Schulz-Baldes - “Edge current channels and Chern numbers in the integer quantum Hall effect” , Rev. Math. Phys. 14 (2002) no. 1, p. 87-119
\bibitem{K16} P. Kuchment - “An overview of periodic elliptic operators” , Bull. Amer. Math. Soc. (N.S.) 53 (2016) no. 3, p. 343-414
\bibitem{L98} P. D. Lax - “On the discriminant of real symmetric matrices” , Comm. Pure Appl. Math. 51 (1998) no. 11-12, p. 1387-1396
\bibitem{L18} M. Lee - “Dirac cones for point scatterers on a honeycomb lattice” , SIAM J. Math. Anal. 48 (2016) no. 2, p. 1459-1488
\bibitem{LWZ17} J. P. Lee-Thorp, M. I. Weinstein \& Y. Zhu - “Elliptic operators with honeycomb symmetry: Dirac points, edge states and applications to photonic graphene” , Arch. Rational Mech. Anal. 232 (2019) no. 1, p. 1-63
\bibitem{M95} I. G. Macdonald - Symmetric functions and Hall polynomials , Oxford Classic Texts in the Physical Sciences , The Clarendon Press, Oxford University Press, New York, 2015
\bibitem{M17} D. Monaco - “Chern and Fu-Kane-Mele invariants as topological obstructions” , in Advances in quantum mechanics , Springer INdAM Ser. , vol. 18 , Springer, Cham, 2017, p. 201-222
\bibitem{M01} J. D. Moore - Lectures on Seiberg-Witten invariants , Lect. Notes in Math. , vol. 1629 , Springer-Verlag, Berlin, 2001
\bibitem{MP14} D. Monaco \& G. Panati - “Topological invariants of eigenvalue intersections and decrease of Wannier functions in graphene” , J. Statist. Phys. 155 (2014) no. 6, p. 1027-1071
\bibitem{P07} G. Panati - “Triviality of Bloch and Bloch-Dirac bundles” , Ann. Inst. H. Poincaré Phys. Théor. 8 (2007) no. 5, p. 995-1011
\bibitem{P02} B. N. Parlett - “The (matrix) discriminant as a determinant” , Linear Algebra Appl. 355 (2002), p. 85-101
\bibitem{P06} P. Petersen - Riemannian geometry , Graduate Texts in Math. , vol. 171 , Springer, Cham, 2016
\bibitem{PS16} E. Prodan \& H. Schulz-Baldes - Bulk and boundary invariants for complex topological insulators , Mathematical Physics Studies , Springer, Cham, 2016
\bibitem{PST02} G. Panati, H. Spohn \& S. Teufel - “Effective dynamics for Bloch electrons: Peierls substitution and beyond” , Comm. Math. Phys. 242 (2003) no. 3, p. 547-578
\bibitem{RH08} S. Raghu \& F. D. M. Haldane - “Analogs of quantum-Hall-effect edge states in photonic crystals” , Phys. Rev. A 78 (2008), article ID 033834, 21 pages
\bibitem{S65} R. T. Seeley - “Extension of $C^{\infty }$ functions defined in a half space” , Proc. Amer. Math. Soc. 15 (1964), p. 625-626
\bibitem{S10} D. Serre - Matrices. Theory and applications , Graduate Texts in Math. , vol. 216 , Springer, New York, 2010
\bibitem{S83} B. Simon - “Holonomy, the quantum adiabatic theorem, and Berry’s phase” , Phys. Rev. Lett. 51 (1983) no. 24, p. 2167-2170
\bibitem{S05} S. F. Singer - Linearity, symmetry, and prediction in the hydrogen atom , Undergraduate Texts in Math. , Springer, New York, 2005
\bibitem{T99} M. Teytel - “How rare are multiple eigenvalues?” , Comm. Pure Appl. Math. 52 (1999) no. 8, p. 917-934
\bibitem{NW29} J. von Neumann \& E. Wigner - “Über das Verhalten von Eigenwerten bei adiabatischen Prozessen” , Phys. Z. 30 (1929), p. 467-470
\bibitem{W47} P. R. Wallace - “The band theory of graphite” , Phys. Rev., II. Ser. 71 (1947), p. 622-634
\end{thebibliography}
\end{document}
